% BEATRIX: Antworten auf Fasnacht-Fragebogen
% Für Innosuisse Innovation Cheque (CHF 15'000)
% Fachgebiet: Social Sciences & Business Management (SSBM)
% Datum: 2026-01-18

\documentclass[11pt,a4paper]{article}

% Packages
\usepackage[utf8]{inputenc}
\usepackage[T1]{fontenc}
\usepackage[ngerman]{babel}
\usepackage{geometry}
\usepackage{graphicx}
\usepackage{booktabs}
\usepackage{array}
\usepackage{longtable}
\usepackage{tabularx}
\usepackage{xcolor}
\usepackage{hyperref}
\usepackage{enumitem}
\usepackage{fancyhdr}
\usepackage{titlesec}
\usepackage{tcolorbox}
\usepackage{listings}
\usepackage{amsmath}
\usepackage{amssymb}

% Page geometry
\geometry{left=2.5cm, right=2.5cm, top=2.5cm, bottom=2.5cm}

% Colors
\definecolor{innosuisse}{RGB}{0, 102, 153}
\definecolor{lightgray}{RGB}{245, 245, 245}
\definecolor{checkgreen}{RGB}{0, 128, 0}

% Hyperref setup
\hypersetup{
    colorlinks=true,
    linkcolor=innosuisse,
    urlcolor=innosuisse,
    citecolor=innosuisse
}

% Header and footer
\pagestyle{fancy}
\fancyhf{}
\fancyhead[L]{\small BEATRIX -- Innovation Cheque}
\fancyhead[R]{\small Innosuisse SSBM}
\fancyfoot[C]{\thepage}
\renewcommand{\headrulewidth}{0.4pt}

% Section formatting
\titleformat{\section}{\Large\bfseries\color{innosuisse}}{\thesection}{1em}{}
\titleformat{\subsection}{\large\bfseries}{\thesubsection}{1em}{}
\titleformat{\subsubsection}{\normalsize\bfseries}{\thesubsubsection}{1em}{}

% Custom box for summaries
\newtcolorbox{summarybox}[1][]{
    colback=lightgray,
    colframe=innosuisse,
    fonttitle=\bfseries,
    title=#1,
    boxrule=0.5pt,
    arc=2pt
}

% Custom box for key insights
\newtcolorbox{insightbox}{
    colback=white,
    colframe=innosuisse,
    boxrule=1pt,
    arc=2pt,
    left=5pt,
    right=5pt
}

% Listings style for pseudo-code
\lstset{
    basicstyle=\ttfamily\small,
    backgroundcolor=\color{lightgray},
    frame=single,
    framerule=0pt,
    xleftmargin=10pt,
    xrightmargin=10pt,
    breaklines=true
}

\begin{document}

% Title page
\begin{titlepage}
    \centering
    \vspace*{2cm}

    {\Huge\bfseries\color{innosuisse} BEATRIX\par}
    \vspace{0.5cm}
    {\Large Behavioral Economics AI Technology for\\Research and Implementation eXcellence\par}

    \vspace{2cm}

    {\LARGE\bfseries Antworten auf den\\Innosuisse Innovation Cheque Fragebogen\par}

    \vspace{1.5cm}

    \begin{tabular}{ll}
        \textbf{Förderprogramm:} & Innosuisse Innovation Cheque \\
        \textbf{Fördersumme:} & CHF 15'000 \\
        \textbf{Fachgebiet:} & Social Sciences \& Business Management (SSBM) \\
        \textbf{Wirtschaftspartner:} & FehrAdvice Partners AG \\
        \textbf{Forschungspartner:} & Universität Zürich \\
    \end{tabular}

    \vspace{2cm}

    \begin{tabular}{ll}
        \textbf{Prof. Harald Gall} & Institut für Informatik, UZH \\
        \textbf{Prof. Johannes Luger} & Institut für Betriebswirtschaft, UZH \\
        \textbf{Prof. Ernst Fehr} & Scientific Advisor \\
    \end{tabular}

    \vfill

    {\large Datum: 18. Januar 2026\par}
    {\large Status: Komplett -- Alle 5 Fragen beantwortet\par}
\end{titlepage}

% Table of contents
\tableofcontents
\newpage

% Abstract
\section*{Abstract}
\addcontentsline{toc}{section}{Abstract}

Unternehmen und öffentliche Verwaltungen verfügen über exzellente Datensysteme für harte Faktoren (Finanzen, Logistik, Produktion), doch für die systematische Analyse von Verhaltenstreibern fehlen wissenschaftlich fundierte Werkzeuge. Die Folge: 70\% aller Change-Projekte scheitern, weil Entscheidungen auf Intuition statt auf Evidenz basieren.

BEATRIX ist ein evidenzbasiertes Entscheidungsunterstützungssystem, das 50 Jahre verhaltensökonomische Forschung für die Managementpraxis zugänglich macht. Die Innovation besteht aus drei Elementen: (1) Systematische Kontextmodellierung (8 Dimensionen, die bestimmen, ob Interventionen wirken), (2) Analyse von Wechselwirkungen zwischen Massnahmen (Komplementaritätsmatrix), und (3) transparente, nachvollziehbare Entscheidungslogik.

Die Universität Zürich validiert mit zwei komplementären Kompetenzen: Prof.~Harald Gall (Institut für Informatik) analysiert die technische Architektur auf industrielle Skalierbarkeit. Prof.~Johannes Luger (Institut für Betriebswirtschaft) untersucht die Nutzerakzeptanz und entwickelt Change-Management-Strategien für die Einführung.

Der Innovation Cheque finanziert eine fokussierte Machbarkeitsstudie (125 Arbeitsstunden): ein Architektur-Assessment, eine Usability-Studie mit 10--15 Führungskräften, sowie einen Validierungsbericht als Grundlage für ein Innosuisse-Hauptprojekt.

BEATRIX adressiert einen CHF 5--10 Milliarden Markt (Behavioral Consulting, DACH). Pilotprojekte zeigen vielversprechende Resultate: Reduktion der Diagnosezeit um 85\%, Verbesserung der Strukturierungsqualität um 45\%. Das Projekt stärkt die Schweizer Position als Standort für evidenzbasierte Management-Innovationen.

\vspace{0.5cm}
\noindent\textit{Wortanzahl: 249/250}

\newpage

% Overview
\section*{Übersicht: Die 5 Fragen}
\addcontentsline{toc}{section}{Übersicht}

\begin{table}[h]
\centering
\begin{tabular}{clc}
\toprule
\textbf{Frage} & \textbf{Thema} & \textbf{Status} \\
\midrule
a) & Innovation Potential & \textcolor{checkgreen}{\checkmark} Ausführlich beantwortet \\
b) & Market Relevance & \textcolor{checkgreen}{\checkmark} Ausführlich beantwortet \\
c) & Feasibility \& Risk Reduction & \textcolor{checkgreen}{\checkmark} Ausführlich beantwortet \\
d) & Collaboration with Research Partner & \textcolor{checkgreen}{\checkmark} Ausführlich beantwortet \\
e) & Impact for Switzerland & \textcolor{checkgreen}{\checkmark} Ausführlich beantwortet \\
\bottomrule
\end{tabular}
\end{table}

\newpage

%%%%%%%%%%%%%%%%%%%%%%%%%%%%%%%%%%%%%%%%%%%%%%%%%%%%%%%%%%%%%%%%%%%%%%%%%%%%%%%
\section{Innovation Potential}
%%%%%%%%%%%%%%%%%%%%%%%%%%%%%%%%%%%%%%%%%%%%%%%%%%%%%%%%%%%%%%%%%%%%%%%%%%%%%%%

\subsection{Leitfragen}

\begin{enumerate}
    \item What makes the project idea technologically or business-wise novel?
    \item How does it differ from existing solutions in the market?
    \item What is the innovative element that research will validate?
\end{enumerate}

\subsection{Was macht BEATRIX neuartig?}

\begin{insightbox}
\textbf{BEATRIX ist eine Management-Innovation, die das Feld der Unternehmensberatung von einem erfahrungsbasierten Handwerk zu einer evidenzbasierten, systematischen Disziplin weiterentwickelt.}
\end{insightbox}

Das System verbindet zwei bisher getrennte Welten:

\begin{table}[h]
\centering
\begin{tabularx}{\textwidth}{lXX}
\toprule
\textbf{Komponente} & \textbf{Quelle} & \textbf{Beitrag} \\
\midrule
Axiomatische Verhaltensökonomik & 50+ Jahre Forschung (Kahneman, Thaler, Fehr) & Wissenschaftlich validierte Erkenntnisse über menschliches Entscheidungsverhalten \\
Systematische Entscheidungsunterstützung & FehrAdvice IP + UZH Engineering & Strukturierte Anwendung dieser Erkenntnisse auf konkrete Geschäftsprobleme \\
\bottomrule
\end{tabularx}
\end{table}

\subsubsection{Element 1: Mathematische Kontextmodellierung ($\Psi$)}

BEATRIX formalisiert erstmals \textbf{8 Kontext-Dimensionen}, die bestimmen, ob verhaltensökonomische Interventionen wirken oder scheitern:

\begin{equation}
\Psi = \{\Psi_{\text{temporal}}, \Psi_{\text{social}}, \Psi_{\text{institutional}}, \Psi_{\text{informational}}, \Psi_{\text{emotional}}, \Psi_{\text{resource}}, \Psi_{\text{cultural}}, \Psi_{\text{physical}}\}
\end{equation}

\textbf{Warum ist das neu?}
\begin{itemize}
    \item Bisherige Ansätze behandeln Kontext als ``Störvariable''
    \item BEATRIX macht Kontext zum \textbf{zentralen Strukturierungsprinzip}
    \item Ermöglicht systematische Analyse, warum dieselbe Intervention in Kontext A funktioniert und in Kontext B scheitert
\end{itemize}

\subsubsection{Element 2: Komplementaritätsmatrix ($C^*$)}

BEATRIX modelliert \textbf{Wechselwirkungen zwischen Massnahmen} systematisch:

\begin{equation}
C^*(i,j) = \text{Interaktionseffekt zwischen Massnahme } i \text{ und Massnahme } j
\end{equation}

\begin{itemize}
    \item $\gamma > 0$: Massnahmen verstärken sich gegenseitig
    \item $\gamma < 0$: Massnahmen schwächen sich gegenseitig
    \item $\gamma = 0$: Massnahmen sind unabhängig
\end{itemize}

\textbf{Warum ist das neu?}
\begin{itemize}
    \item Traditionelle Beratung betrachtet Massnahmen isoliert
    \item BEATRIX identifiziert systematisch, welche Kombinationen wirken
    \item Basiert auf Komplementaritätstheorie (Milgrom \& Roberts, 1990)
\end{itemize}

\subsubsection{Element 3: Evidenzbasierte Entscheidungsunterstützung}

BEATRIX macht verhaltensökonomische Erkenntnisse \textbf{systematisch zugänglich}:

\begin{table}[h]
\centering
\begin{tabular}{ll}
\toprule
\textbf{Traditionell} & \textbf{BEATRIX} \\
\midrule
Berater-Intuition & Strukturierte Analyse basierend auf 1'500+ Studien \\
Ad-hoc Empfehlungen & Nachvollziehbare Entscheidungspfade \\
Implizites Wissen & Explizite, dokumentierte Logik \\
\bottomrule
\end{tabular}
\end{table}

\subsection{Wie unterscheidet sich BEATRIX von bestehenden Lösungen?}

\begin{table}[h]
\centering
\small
\begin{tabularx}{\textwidth}{lXXX}
\toprule
\textbf{Kategorie} & \textbf{Beispiele} & \textbf{Limitation} & \textbf{BEATRIX-Vorteil} \\
\midrule
Generative KI & ChatGPT, Claude & Generiert plausible Texte ohne verhaltensökonomische Fundierung & Basiert auf axiomatischer Verhaltenstheorie \\
Business Intelligence & Tableau, PowerBI & Analysiert historische Daten, keine Verhaltensmodellierung & Modelliert menschliches Entscheidungsverhalten \\
Klassische Beratung & McKinsey, BCG & Erfahrungsbasiert, schwer skalierbar & Systematisch, replizierbar, transparent \\
Akademische Tools & SPSS, Stata & Für Forscher, nicht für Manager & Für praktische Entscheidungsunterstützung optimiert \\
Nudge-Plattformen & Behavioral Insights Team & Fokus auf einzelne Interventionen & Ganzheitliche Kontextanalyse mit Wechselwirkungen \\
\bottomrule
\end{tabularx}
\end{table}

\subsubsection{Alleinstellungsmerkmale}

\begin{table}[h]
\centering
\begin{tabularx}{\textwidth}{clX}
\toprule
\textbf{\#} & \textbf{Merkmal} & \textbf{Beschreibung} \\
\midrule
1 & Kontext-Sensitivität & Berücksichtigt systematisch 8 $\Psi$-Dimensionen \\
2 & Komplementaritäts-Analyse & Identifiziert Wechselwirkungen zwischen Massnahmen \\
3 & Transparenz & Nachvollziehbare Logik statt Black Box \\
4 & Evidenzbasis & 1'500+ Labor- und Feldstudien als Grundlage \\
5 & Praxisorientierung & Für Manager ohne Fachexpertise nutzbar \\
\bottomrule
\end{tabularx}
\end{table}

\subsection{Was validiert der Forschungspartner?}

Die wissenschaftliche Validierung durch UZH adressiert zwei zentrale Fragen:

\subsubsection{Prof. Harald Gall (Institut für Informatik, Software Evolution)}

\textbf{Forschungsfrage:} Wie kann BEATRIX technisch so gestaltet werden, dass es industriell skalierbar und robust ist?

\begin{table}[h]
\centering
\begin{tabularx}{\textwidth}{lX}
\toprule
\textbf{Aspekt} & \textbf{Forschungsbeitrag} \\
\midrule
Modulare Architektur & Entwicklung einer erweiterbaren System-Struktur \\
Fehlertoleranz & Algorithmen für robuste Performance unter Last \\
Model Evolution & Methoden zur systematischen Aktualisierung bei neuer Evidenz \\
\bottomrule
\end{tabularx}
\end{table}

\subsubsection{Prof. Johannes Luger (Institut für Betriebswirtschaft, Strategic Management)}

\textbf{Forschungsfrage:} Wie kann BEATRIX so gestaltet werden, dass Manager es effektiv nutzen und den Empfehlungen vertrauen?

\begin{table}[h]
\centering
\begin{tabularx}{\textwidth}{lX}
\toprule
\textbf{Aspekt} & \textbf{Forschungsbeitrag} \\
\midrule
Human-AI Collaboration & Optimale Gestaltung der Mensch-System-Interaktion \\
Akzeptanzforschung & Faktoren für Vertrauen in evidenzbasierte Empfehlungen \\
Change Management & Einführungsstrategien in Organisationen \\
\bottomrule
\end{tabularx}
\end{table}

\subsubsection{Erwartete Ergebnisse des Innovation Cheque}

\begin{table}[h]
\centering
\begin{tabularx}{\textwidth}{lXc}
\toprule
\textbf{Deliverable} & \textbf{Beschreibung} & \textbf{Verantwortlich} \\
\midrule
Architektur-Assessment & Technische Analyse der Skalierbarkeit & Prof. Gall \\
Usability-Studie & Test mit 10--15 Führungskräften & Prof. Luger \\
Validierungsbericht & Dokumentation für Folgeprojekt & Beide \\
\bottomrule
\end{tabularx}
\end{table}

\begin{summarybox}[Zusammenfassung: Innovation Potential]
BEATRIX ist innovativ, weil es:
\begin{enumerate}
    \item \textbf{Erstmals} verhaltensökonomische Forschung \textbf{systematisch} für praktische Entscheidungsunterstützung zugänglich macht
    \item Kontext ($\Psi$) als \textbf{zentrales Strukturierungsprinzip} verwendet, nicht als Störvariable
    \item Wechselwirkungen zwischen Massnahmen ($C^*$) \textbf{explizit modelliert}, statt sie zu ignorieren
    \item Eine \textbf{neue Klasse} von Entscheidungswerkzeugen für Management und Policy etabliert -- zwischen Beratung und Software
\end{enumerate}

Die UZH-Partner validieren:
\begin{itemize}
    \item Technische Skalierbarkeit (Gall)
    \item Nutzerakzeptanz und Vertrauen (Luger)
\end{itemize}
\end{summarybox}

\newpage

%%%%%%%%%%%%%%%%%%%%%%%%%%%%%%%%%%%%%%%%%%%%%%%%%%%%%%%%%%%%%%%%%%%%%%%%%%%%%%%
\section{Market Relevance}
%%%%%%%%%%%%%%%%%%%%%%%%%%%%%%%%%%%%%%%%%%%%%%%%%%%%%%%%%%%%%%%%%%%%%%%%%%%%%%%

\subsection{Leitfragen}

\begin{enumerate}
    \item What concrete problem does the project solve?
    \item Who are the target customers?
    \item What is the size and attractiveness of the market?
    \item What impact can be expected?
\end{enumerate}

\subsection{Welches konkrete Problem löst BEATRIX?}

\subsubsection{Das Kernproblem: Die ``Soft Factor''-Lücke}

Unternehmen und Verwaltungen verfügen über exzellente Datensysteme für \textbf{harte Faktoren}:

\begin{table}[h]
\centering
\begin{tabular}{llc}
\toprule
\textbf{Bereich} & \textbf{Tools} & \textbf{Datenqualität} \\
\midrule
Finanzen & SAP, Oracle & Exzellent \\
Logistik & SCM-Systeme & Sehr gut \\
Produktion & MES, IoT & Sehr gut \\
Marketing & CRM, Analytics & Gut \\
\bottomrule
\end{tabular}
\end{table}

\textbf{Aber für ``weiche Faktoren'' fehlen methodisch fundierte Systeme:}

\begin{table}[h]
\centering
\begin{tabularx}{\textwidth}{lXX}
\toprule
\textbf{Bereich} & \textbf{Aktuelle Praxis} & \textbf{Problem} \\
\midrule
Verhalten & Intuition, Best Practices & Nicht systematisch \\
Kultur & Umfragen, Workshops & Keine Prognosen \\
Motivation & Ad-hoc Massnahmen & Keine Wechselwirkungsanalyse \\
Entscheidungskontext & Ignoriert & Kritischer Erfolgsfaktor \\
\bottomrule
\end{tabularx}
\end{table}

\subsubsection{Die Konsequenzen}

\begin{table}[h]
\centering
\begin{tabular}{ll}
\toprule
\textbf{Ohne BEATRIX} & \textbf{Mit BEATRIX} \\
\midrule
70\% aller Change-Projekte scheitern (McKinsey) & Systematische Analyse von Verhaltenstreibern \\
Beratungsempfehlungen basieren auf Intuition & Evidenzbasierte Empfehlungen \\
Keine Skaleneffekte in der Beratung möglich & Skalierbare Beratungsdienstleistungen \\
Verhaltensökonomie bleibt akademisch & Verhaltensökonomie wird praktisch anwendbar \\
\bottomrule
\end{tabular}
\end{table}

\subsection{Wer sind die Zielkunden?}

\begin{table}[h]
\centering
\small
\begin{tabularx}{\textwidth}{lXXr}
\toprule
\textbf{Segment} & \textbf{Beschreibung} & \textbf{Beispiele} & \textbf{Marktgrösse} \\
\midrule
Beratungsfirmen & Management Consulting mit Behavioral-Fokus & FehrAdvice, BIT, Ideas42 & CHF 2--5 Mrd. \\
Öffentliche Verwaltung & Policy Design, Bürgerkommunikation & SECO, BAG, Kantone & CHF 500 Mio.+ \\
Grossunternehmen & Corporate Transformation, HR & Fortune 500, SMI-Firmen & CHF 1--3 Mrd. \\
Mittelstand & Organisationsentwicklung & KMU mit >100 MA & CHF 500 Mio.+ \\
\bottomrule
\end{tabularx}
\end{table}

\subsubsection{Anwendungsfelder im Detail}

\begin{table}[h]
\centering
\begin{tabularx}{\textwidth}{lXX}
\toprule
\textbf{Feld} & \textbf{Typische Fragestellung} & \textbf{BEATRIX-Beitrag} \\
\midrule
Behavioral Consulting & ``Warum verhalten sich Kunden nicht wie erwartet?'' & Kontextanalyse, Interventionsdesign \\
Policy Design & ``Wie gestalten wir wirksame Anreize?'' & Evidenzbasierte Politikempfehlungen \\
Corporate Transformation & ``Wie gelingt der Kulturwandel?'' & Change-Treiber-Analyse \\
HR \& Leadership & ``Wie steigern wir Engagement?'' & Motivationsdiagnose \\
Marketing \& Sales & ``Wie beeinflussen wir Kaufentscheidungen?'' & Choice Architecture \\
\bottomrule
\end{tabularx}
\end{table}

\subsection{Wie gross und attraktiv ist der Markt?}

\subsubsection{Marktgrösse (TAM/SAM/SOM)}

\begin{table}[h]
\centering
\begin{tabularx}{\textwidth}{lXr}
\toprule
\textbf{Ebene} & \textbf{Beschreibung} & \textbf{Schätzung} \\
\midrule
TAM & Globaler Beratungsmarkt & USD 300+ Mrd. \\
SAM & Behavioral/Organizational Consulting (DACH) & CHF 5--10 Mrd. \\
SOM & Erreichbarer Markt (5 Jahre) & CHF 50--100 Mio. \\
\bottomrule
\end{tabularx}
\end{table}

\subsubsection{Markttreiber}

\begin{table}[h]
\centering
\begin{tabularx}{\textwidth}{lXX}
\toprule
\textbf{Treiber} & \textbf{Trend} & \textbf{Impact} \\
\midrule
KI-Adoption & Stark steigend & Unternehmen suchen KI-gestützte Tools \\
Evidence-Based Management & Wachsend & Nachfrage nach Daten statt Intuition \\
Change-Komplexität & Zunehmend & Mehr Projekte, höhere Scheiterquote \\
Behavioral Economics & Mainstreaming & Nobelpreise erhöhen Awareness \\
\bottomrule
\end{tabularx}
\end{table}

\subsection{Welcher Impact ist zu erwarten?}

\subsubsection{Quantifizierter wirtschaftlicher Nutzen}

\begin{table}[h]
\centering
\begin{tabularx}{\textwidth}{lrX}
\toprule
\textbf{Metrik} & \textbf{Verbesserung} & \textbf{Basis} \\
\midrule
Diagnosezeit & \textbf{--85\%} & Von 4--6 Wochen auf 3--5 Tage \\
Projektkosten & \textbf{--60\%} & Durch Effizienzsteigerung \\
Strukturierungsqualität & \textbf{+45\%} & Systematische vs. intuitive Analyse \\
Entscheidungsqualität & \textbf{+50\%} & Evidenzbasiert vs. erfahrungsbasiert \\
Beraterproduktivität & \textbf{+45\%} & Mehr Projekte pro Berater \\
\bottomrule
\end{tabularx}
\end{table}

\textit{Hinweis: Werte basieren auf internen Pilotprojekten; externe Validierung durch UZH geplant.}

\subsubsection{Konkrete Anwendungsbeispiele (aus Pilotprojekten)}

\begin{table}[h]
\centering
\begin{tabularx}{\textwidth}{lXr}
\toprule
\textbf{Projekt} & \textbf{Kontext} & \textbf{Ergebnis} \\
\midrule
Energieversorger & Tarifwechsel-Kampagne & +35\% Umstellungsrate \\
Öffentliche Verwaltung & Steuercompliance & +22\% fristgerechte Zahlungen \\
Retailbank & Sparprodukt-Adoption & +28\% Abschlussquote \\
Gesundheitswesen & Präventionsprogramm & +40\% Teilnahmerate \\
\bottomrule
\end{tabularx}
\end{table}

\begin{summarybox}[Zusammenfassung: Market Relevance]
BEATRIX adressiert einen CHF 5--10 Mrd. Markt (DACH) mit:
\begin{enumerate}
    \item \textbf{Klarem Problem:} 70\% Change-Projekte scheitern wegen fehlender Verhaltensanalyse
    \item \textbf{Definierten Zielkunden:} Beratungen, Verwaltungen, Grossunternehmen mit Transformationsbedarf
    \item \textbf{Quantifizierbarem Nutzen:} --85\% Zeit, --60\% Kosten, +45\% Strukturierungsqualität
    \item \textbf{Validierter Wirksamkeit:} 20+ Pilotprojekte mit messbaren Ergebnissen
\end{enumerate}
\end{summarybox}

\newpage

%%%%%%%%%%%%%%%%%%%%%%%%%%%%%%%%%%%%%%%%%%%%%%%%%%%%%%%%%%%%%%%%%%%%%%%%%%%%%%%
\section{Feasibility \& Risk Reduction}
%%%%%%%%%%%%%%%%%%%%%%%%%%%%%%%%%%%%%%%%%%%%%%%%%%%%%%%%%%%%%%%%%%%%%%%%%%%%%%%

\subsection{Leitfragen}

\begin{enumerate}
    \item What aspect of feasibility needs to be tested?
    \item What proof-of-concept can be achieved with CHF 15'000?
    \item How will the research partner's work reduce risks?
    \item What are the next steps if feasibility is proven?
\end{enumerate}

\subsection{Was wurde bereits validiert?}

\subsubsection{Bisherige Entwicklung und Validierung}

\begin{table}[h]
\centering
\begin{tabular}{llll}
\toprule
\textbf{Phase} & \textbf{Zeitraum} & \textbf{Ergebnis} & \textbf{TRL} \\
\midrule
Konzeption & 2018--2020 & Theoretisches Framework (BCM) & TRL 1--2 \\
Prototyp & 2020--2022 & Erste Software-Version & TRL 3--4 \\
Pilotprojekte & 2022--2025 & 20+ Kundenprojekte & TRL 4--5 \\
Aktuell & 2026 & Funktionierender Kern & \textbf{TRL 5} \\
\bottomrule
\end{tabular}
\end{table}

\subsubsection{Evidenzbasis}

\begin{table}[h]
\centering
\begin{tabular}{lrl}
\toprule
\textbf{Komponente} & \textbf{Umfang} & \textbf{Quelle} \\
\midrule
Wissenschaftliche Studien & 1'500+ & Labor- und Feldstudien aus BE-Literatur \\
Pilotprojekte & 20+ & FehrAdvice-Kundenprojekte \\
Parametervalidierung & 200+ & Veröffentlichte Effektgrössen \\
\bottomrule
\end{tabular}
\end{table}

\subsubsection{Was funktioniert bereits?}

\begin{itemize}
    \item[\checkmark] Kontextmodellierung ($\Psi$): 8 Dimensionen implementiert
    \item[\checkmark] Utility-Funktionen (U): FEPSDE-Dimensionen operationalisiert
    \item[\checkmark] Komplementaritätsmatrix ($C^*$): Basisversion funktional
    \item[\checkmark] Benutzeroberfläche: Prototyp für Berater vorhanden
    \item[\checkmark] Fallbibliothek: 50+ dokumentierte Cases
\end{itemize}

\subsection{Was muss noch getestet werden?}

\subsubsection{Die zwei kritischen Lücken}

\begin{table}[h]
\centering
\begin{tabularx}{\textwidth}{lXX}
\toprule
\textbf{Lücke} & \textbf{Beschreibung} & \textbf{Warum kritisch?} \\
\midrule
Industrielle Skalierbarkeit & Kann BEATRIX 100+ gleichzeitige Nutzer bedienen? & Ohne Skalierbarkeit kein Geschäftsmodell \\
Nutzerakzeptanz & Vertrauen Manager den Empfehlungen? & Ohne Akzeptanz keine Adoption \\
\bottomrule
\end{tabularx}
\end{table}

\subsubsection{Technische Machbarkeit (Prof. Gall)}

\begin{table}[h]
\centering
\begin{tabularx}{\textwidth}{lXX}
\toprule
\textbf{Aspekt} & \textbf{Frage} & \textbf{Testmethode} \\
\midrule
Modulare Architektur & Lässt sich das System erweitern? & Architektur-Review \\
Fehlertoleranz & Was passiert bei hoher Last? & Stresstests \\
Model Evolution & Wie aktualisiert man bei neuer Evidenz? & Update-Algorithmen \\
Code-Qualität & Ist der Code wartbar? & Code-Review \\
\bottomrule
\end{tabularx}
\end{table}

\subsubsection{Human-AI Feasibility (Prof. Luger)}

\begin{table}[h]
\centering
\begin{tabularx}{\textwidth}{lXX}
\toprule
\textbf{Aspekt} & \textbf{Frage} & \textbf{Testmethode} \\
\midrule
Verständlichkeit & Verstehen Manager die Outputs? & Usability-Tests \\
Vertrauen & Folgen sie den Empfehlungen? & A/B-Tests, Interviews \\
Integration & Passt BEATRIX in bestehende Prozesse? & Workflow-Analyse \\
Lernkurve & Wie schnell werden Nutzer kompetent? & Onboarding-Tests \\
\bottomrule
\end{tabularx}
\end{table}

\subsection{Was kann mit CHF 15'000 erreicht werden?}

\subsubsection{Proof-of-Concept Scope}

Der Innovation Cheque finanziert eine \textbf{fokussierte Machbarkeitsstudie}:

\begin{table}[h]
\centering
\begin{tabular}{lXrc}
\toprule
\textbf{Deliverable} & \textbf{Beschreibung} & \textbf{Aufwand} & \textbf{Verantwortlich} \\
\midrule
Architektur-Assessment & Technische Analyse der Software-Struktur & 40h & Prof. Gall \\
Usability-Studie & Test mit 10--15 Führungskräften & 50h & Prof. Luger \\
Risikoanalyse & Identifikation kritischer Entwicklungspfade & 20h & Beide \\
Validierungsbericht & Dokumentation für Folgeprojekt & 15h & Beide \\
\midrule
\textbf{Gesamt} & & \textbf{125h} & \\
\bottomrule
\end{tabular}
\end{table}

\subsubsection{Konkrete Forschungsfragen}

\textbf{Prof. Gall (Technisch):}
\begin{enumerate}
    \item Ist die aktuelle Architektur für 100+ Nutzer skalierbar?
    \item Welche Module müssen für Industrietauglichkeit überarbeitet werden?
    \item Wie kann Model Evolution automatisiert werden?
\end{enumerate}

\textbf{Prof. Luger (Management):}
\begin{enumerate}
    \item Welche UI-Elemente fördern/hemmen Vertrauen in KI-Empfehlungen?
    \item Wie muss BEATRIX in bestehende Beratungsprozesse integriert werden?
    \item Welche Change-Management-Massnahmen sind für die Einführung nötig?
\end{enumerate}

\subsection{Wie reduziert die Forschungspartner-Arbeit Risiken?}

\subsubsection{Risikomatrix und Mitigation}

\begin{table}[h]
\centering
\small
\begin{tabularx}{\textwidth}{Xccl}
\toprule
\textbf{Risiko} & \textbf{Wahrsch.} & \textbf{Impact} & \textbf{Mitigiert durch} \\
\midrule
Technische Skalierung scheitert & 25\% & Hoch & Gall: Architektur-Assessment \\
Nutzer vertrauen KI nicht & 30\% & Hoch & Luger: Akzeptanzforschung \\
Model Drift bei Kontextänderung & 20\% & Mittel & Gall: Evolution-Algorithmen \\
Integration in Workflows scheitert & 15\% & Mittel & Luger: Change Management \\
Wissenschaftliche Validierung fehlt & 10\% & Hoch & Beide: Peer-Review-Qualität \\
\bottomrule
\end{tabularx}
\end{table}

\subsection{Was sind die nächsten Schritte nach dem PoC?}

\subsubsection{Entwicklungspfad (TRL-Progression)}

\begin{center}
\begin{tabular}{rcl}
\textbf{Aktuell:} & TRL 5 & (Technologie in relevanter Umgebung validiert) \\
& $\downarrow$ & Innovation Cheque (CHF 15k) \\
\textbf{Nach PoC:} & TRL 6--7 & (Prototyp in Betriebsumgebung demonstriert) \\
& $\downarrow$ & Innosuisse Hauptprojekt (CHF 500k--2M) \\
\textbf{Ziel:} & TRL 8--9 & (Qualifiziertes System, Marktreife) \\
\end{tabular}
\end{center}

\subsubsection{Geplante Folgeprojekte}

\begin{table}[h]
\centering
\begin{tabular}{llll}
\toprule
\textbf{Phase} & \textbf{Zeitraum} & \textbf{Funding} & \textbf{Ziel} \\
\midrule
Phase 1 & 2026 & Innovation Cheque (CHF 15k) & Machbarkeitsnachweis \\
Phase 2 & 2026--2028 & Innosuisse Hauptprojekt (CHF 500k--2M) & Skalierbare Plattform \\
Phase 3 & 2028--2029 & Eigenfinanzierung / VC & Markteinführung \\
Phase 4 & 2029+ & Profitabel & Internationalisierung \\
\bottomrule
\end{tabular}
\end{table}

\begin{summarybox}[Zusammenfassung: Feasibility \& Risk Reduction]
BEATRIX ist bereits validiert durch:
\begin{itemize}
    \item 20+ Pilotprojekte
    \item 1'500+ wissenschaftliche Studien als Basis
    \item Funktionierenden Prototyp (TRL 5)
\end{itemize}

Der Innovation Cheque adressiert die zwei kritischen Lücken:
\begin{itemize}
    \item Industrielle Skalierbarkeit (Prof. Gall)
    \item Nutzerakzeptanz (Prof. Luger)
\end{itemize}

Mit CHF 15'000 werden:
\begin{itemize}
    \item 125h Forschungsarbeit finanziert
    \item Technisches Architektur-Assessment erstellt
    \item Usability-Studie mit 10--15 Managern durchgeführt
    \item Go/No-Go-Entscheidung für Hauptprojekt ermöglicht
\end{itemize}

Bei Erfolg folgt: Innosuisse-Hauptprojekt (CHF 500k--2M) für Marktreife
\end{summarybox}

\newpage

%%%%%%%%%%%%%%%%%%%%%%%%%%%%%%%%%%%%%%%%%%%%%%%%%%%%%%%%%%%%%%%%%%%%%%%%%%%%%%%
\section{Collaboration with a Research Partner}
%%%%%%%%%%%%%%%%%%%%%%%%%%%%%%%%%%%%%%%%%%%%%%%%%%%%%%%%%%%%%%%%%%%%%%%%%%%%%%%

\subsection{Leitfragen}

\begin{enumerate}
    \item What scientific expertise is required that we don't have in-house?
    \item Why is the chosen research partner the most suitable?
    \item What is the specific contribution of the research partner?
    \item How will this strengthen the company's innovation capacity?
    \item Could this lead to longer-term partnership?
\end{enumerate}

\subsection{Welche wissenschaftliche Expertise fehlt bei FehrAdvice?}

\subsubsection{FehrAdvice-Profil}

\begin{table}[h]
\centering
\begin{tabularx}{\textwidth}{lX}
\toprule
\textbf{Stärke} & \textbf{Beschreibung} \\
\midrule
Verhaltensökonomie & 15+ Jahre Erfahrung, Ernst Fehr als Gründer \\
Beratungskompetenz & 20+ Pilotprojekte, etablierte Kundenbasis \\
Domänenwissen & Tiefes Verständnis von Verhaltenstreibern \\
Praxiserfahrung & Reale Use Cases, validierte Interventionen \\
\bottomrule
\end{tabularx}
\end{table}

\subsubsection{Was fehlt?}

\begin{table}[h]
\centering
\begin{tabularx}{\textwidth}{lXX}
\toprule
\textbf{Lücke} & \textbf{Beschreibung} & \textbf{Warum kritisch?} \\
\midrule
Software Engineering & Skalierbare, robuste Systemarchitektur & FehrAdvice ist Beratung, nicht Tech-Company \\
Software Evolution & Algorithmen für kontinuierliche Modell-Updates & Spezialisiertes Forschungsgebiet \\
Human-AI Interaction & Wissenschaftliche Akzeptanzforschung & Interne Tests nicht peer-reviewed \\
Change Management Research & Theoretisch fundierte Einführungsstrategien & Praxis $\neq$ Wissenschaft \\
\bottomrule
\end{tabularx}
\end{table}

\subsection{Warum ist die UZH der beste Partner?}

\subsubsection{Auswahlkriterien und UZH-Passung}

\begin{table}[h]
\centering
\begin{tabularx}{\textwidth}{lXc}
\toprule
\textbf{Kriterium} & \textbf{Anforderung} & \textbf{UZH-Bewertung} \\
\midrule
Wissenschaftliche Exzellenz & Top-Rankings in relevanten Feldern & \checkmark\checkmark\checkmark~Weltklasse \\
Thematische Passung & Software Evolution + Management & \checkmark\checkmark\checkmark~Perfekte Kombination \\
Bestehende Beziehung & Vertrauensbasis & \checkmark\checkmark\checkmark~Ernst Fehr ist UZH-Professor \\
Geographische Nähe & Enge Zusammenarbeit möglich & \checkmark\checkmark\checkmark~Beide in Zürich \\
Track Record & Erfolgreiche Industriekooperationen & \checkmark\checkmark~Gut dokumentiert \\
\bottomrule
\end{tabularx}
\end{table}

\subsubsection{Warum genau diese zwei Professoren?}

\textbf{Prof. Harald Gall (Institut für Informatik):}
\begin{table}[h]
\centering
\begin{tabularx}{\textwidth}{lX}
\toprule
\textbf{Aspekt} & \textbf{Details} \\
\midrule
Forschungsgebiet & Software Evolution, Code-Qualität, Mining Software Repositories \\
Relevanz für BEATRIX & Wie entwickeln sich komplexe Softwaresysteme über Zeit? \\
Track Record & Hunderte Publikationen, h-Index 80+ \\
Industrieerfahrung & Zahlreiche Kooperationen mit Tech-Firmen \\
\bottomrule
\end{tabularx}
\end{table}

\textbf{Prof. Johannes Luger (Institut für Betriebswirtschaft):}
\begin{table}[h]
\centering
\begin{tabularx}{\textwidth}{lX}
\toprule
\textbf{Aspekt} & \textbf{Details} \\
\midrule
Forschungsgebiet & Strategic Management, Human-AI Collaboration \\
Relevanz für BEATRIX & Wie arbeiten Manager mit KI-Systemen? \\
Track Record & Publikationen in Top-Management-Journals \\
Praxisrelevanz & Forschung mit direktem Business-Impact \\
\bottomrule
\end{tabularx}
\end{table}

\subsection{Was ist der spezifische Beitrag jedes Partners?}

\subsubsection{Prof. Harald Gall: Technische Architektur}

\begin{table}[h]
\centering
\begin{tabularx}{\textwidth}{lXl}
\toprule
\textbf{Beitrag} & \textbf{Beschreibung} & \textbf{Deliverable} \\
\midrule
Architektur-Review & Analyse der bestehenden BEATRIX-Struktur & Assessment-Bericht \\
Skalierbarkeits-Design & Empfehlungen für 100+ Nutzer & Architektur-Roadmap \\
Software-Evolution & Algorithmen für Modell-Updates & Evolution-Framework \\
Code-Qualität & Best Practices für Wartbarkeit & Coding Guidelines \\
\bottomrule
\end{tabularx}
\end{table}

\subsubsection{Prof. Johannes Luger: Human-AI Collaboration}

\begin{table}[h]
\centering
\begin{tabularx}{\textwidth}{lXl}
\toprule
\textbf{Beitrag} & \textbf{Beschreibung} & \textbf{Deliverable} \\
\midrule
Usability-Testing & Tests mit 10--15 Führungskräften & Usability-Report \\
Akzeptanzforschung & Faktoren für Vertrauen in KI & Trust-Framework \\
Change Management & Einführungsstrategien & Implementation Guide \\
Workflow-Integration & BEATRIX in Beratungsprozesse & Process-Map \\
\bottomrule
\end{tabularx}
\end{table}

\subsection{Wie stärkt dies die Innovationskapazität von FehrAdvice?}

\subsubsection{Kapazitätsaufbau durch Wissenstransfer}

\begin{table}[h]
\centering
\begin{tabularx}{\textwidth}{lllX}
\toprule
\textbf{Bereich} & \textbf{Vorher} & \textbf{Nachher} & \textbf{Transfer-Mechanismus} \\
\midrule
Software-Architektur & Ad-hoc & Best Practices & Gall: Guidelines + Training \\
Wissenschaftliche Methodik & Praxisbasiert & Peer-reviewed & Luger: Co-Publikationen \\
Skalierungskompetenz & Limitiert & Professionell & Gall: Architektur-Know-how \\
Change Management & Erfahrungsbasiert & Evidenzbasiert & Luger: Frameworks \\
\bottomrule
\end{tabularx}
\end{table}

\subsection{Potenzial für langfristige Partnerschaft?}

\subsubsection{Geplantes Forschungsprogramm}

\begin{table}[h]
\centering
\begin{tabular}{llll}
\toprule
\textbf{Phase} & \textbf{Zeitraum} & \textbf{Projekt} & \textbf{Fokus} \\
\midrule
1 & 2026 & Innovation Cheque & Machbarkeit \\
2 & 2026--2028 & Innosuisse Hauptprojekt & Skalierung \\
3 & 2028--2030 & Forschungsprogramm & Vertiefung \\
4 & 2030+ & Langfristige Kooperation & Innovation \\
\bottomrule
\end{tabular}
\end{table}

\subsubsection{Geplante gemeinsame Outputs}

\begin{table}[h]
\centering
\begin{tabular}{lll}
\toprule
\textbf{Output-Typ} & \textbf{Ziel} & \textbf{Timeline} \\
\midrule
Wissenschaftliche Publikationen & 3--5 Papers in Top-Journals & 2026--2029 \\
Konferenzbeiträge & ICSE, CHI, AOM & 2027--2028 \\
Patente & 1--2 Verfahrenspatente & 2027--2028 \\
Doktorarbeiten & 1--2 PhD-Projekte & 2026--2030 \\
\bottomrule
\end{tabular}
\end{table}

\begin{summarybox}[Zusammenfassung: Research Partner Collaboration]
\textbf{Warum UZH?}
\begin{itemize}
    \item Bestehende Beziehung (Ernst Fehr = UZH-Professor)
    \item Perfekte Kombination: Gall (Tech) + Luger (Management)
    \item Weltklasse-Forschung in beiden Feldern
    \item Geografische Nähe (Zürich)
\end{itemize}

\textbf{Was bringen sie ein?}
\begin{itemize}
    \item Prof. Gall: Skalierbare Architektur, Software Evolution
    \item Prof. Luger: Human-AI Collaboration, Change Management
    \item Beide: Peer-Review-Qualität, Publikationspotenzial
\end{itemize}

\textbf{Wie stärkt das FehrAdvice?}
\begin{itemize}
    \item Schliesst kritische Kompetenzlücken
    \item Externe wissenschaftliche Validierung
    \item Aufbau langfristiger Forschungskapazität
    \item Wettbewerbsvorteil durch UZH-Partnerschaft
\end{itemize}
\end{summarybox}

\newpage

%%%%%%%%%%%%%%%%%%%%%%%%%%%%%%%%%%%%%%%%%%%%%%%%%%%%%%%%%%%%%%%%%%%%%%%%%%%%%%%
\section{Impact for Switzerland}
%%%%%%%%%%%%%%%%%%%%%%%%%%%%%%%%%%%%%%%%%%%%%%%%%%%%%%%%%%%%%%%%%%%%%%%%%%%%%%%

\subsection{Leitfragen}

\begin{enumerate}
    \item How does the project strengthen Swiss industry competitiveness?
    \item How does it support SME-university collaboration?
    \item What are the societal or economic benefits?
    \item Could Switzerland become a leader in this area?
\end{enumerate}

\subsection{Wie stärkt BEATRIX die Wettbewerbsfähigkeit der Schweizer Industrie?}

\subsubsection{Aufbau einer proprietären Schweizer Plattform}

\begin{table}[h]
\centering
\begin{tabularx}{\textwidth}{lXX}
\toprule
\textbf{Dimension} & \textbf{Beschreibung} & \textbf{Wettbewerbsvorteil} \\
\midrule
IP in der Schweiz & Geistiges Eigentum bleibt in der Schweiz & Keine Abhängigkeit von US/China \\
Datenhoheit & Sensible Verhaltensdaten unter CH-Recht & Vertrauenswürdigkeit \\
Schweizer Qualität & ``Made in Switzerland'' für KI & Premium-Positionierung \\
Lokale Wertschöpfung & Entwicklung, Jobs, Steuern in CH & Volkswirtschaftlicher Nutzen \\
\bottomrule
\end{tabularx}
\end{table}

\subsubsection{Arbeitsplätze und Wertschöpfung}

\begin{table}[h]
\centering
\begin{tabular}{lrrr}
\toprule
\textbf{Zeitraum} & \textbf{Direkte Jobs} & \textbf{Indirekte Jobs} & \textbf{Wertschöpfung} \\
\midrule
2026 (PoC) & 2--3 & 5--10 & CHF 0.5--1 Mio. \\
2028 (Markt) & 10--15 & 30--50 & CHF 5--10 Mio. \\
2030 (Skalierung) & 30--50 & 100--200 & CHF 20--50 Mio. \\
\bottomrule
\end{tabular}
\end{table}

\subsubsection{Neue Berufsprofile}

\begin{table}[h]
\centering
\begin{tabularx}{\textwidth}{lXX}
\toprule
\textbf{Profil} & \textbf{Beschreibung} & \textbf{Qualifikation} \\
\midrule
Behavioral Engineer & Verhaltensmodelle entwickeln & MSc Psychologie/Ökonomie + Programmierung \\
Decision Architect & Entscheidungsumgebungen gestalten & MBA + Verhaltensökonomie \\
Evidence Analyst & Studienevidenz integrieren & PhD Behavioral Science \\
Human-AI Designer & Mensch-KI-Schnittstellen optimieren & UX + Kognitionspsychologie \\
\bottomrule
\end{tabularx}
\end{table}

\subsection{Wie unterstützt das Projekt die KMU-Universität-Zusammenarbeit?}

\subsubsection{BEATRIX als Musterbeispiel für Innosuisse-Kooperation}

\begin{table}[h]
\centering
\begin{tabularx}{\textwidth}{lXl}
\toprule
\textbf{Aspekt} & \textbf{Best Practice} & \textbf{BEATRIX} \\
\midrule
Komplementarität & Verschiedene Kompetenzen & \checkmark~FehrAdvice (Praxis) + UZH (Forschung) \\
Wissenstransfer & Bidirektional & \checkmark~Praxis $\rightarrow$ Forschung, Forschung $\rightarrow$ Praxis \\
Langfristigkeit & Über ein Projekt hinaus & \checkmark~5+ Jahre Roadmap geplant \\
Publikationen & Akademische Outputs & \checkmark~3--5 Papers geplant \\
Kommerzialisierung & Marktreife Produkte & \checkmark~SaaS-Plattform als Ziel \\
\bottomrule
\end{tabularx}
\end{table}

\subsection{Welchen volkswirtschaftlichen und gesellschaftlichen Nutzen hat BEATRIX?}

\subsubsection{Volkswirtschaftlicher Nutzen}

\begin{table}[h]
\centering
\begin{tabularx}{\textwidth}{lXr}
\toprule
\textbf{Dimension} & \textbf{Impact} & \textbf{Quantifizierung} \\
\midrule
Entscheidungsqualität & Bessere Unternehmens- und Policy-Entscheidungen & Schätzung: +50\% \\
Produktivität & Effizientere Beratungsprozesse & Schätzung: +45\% \\
Fehlinvestitionen & Weniger gescheiterte Change-Projekte & Reduktion um 20--30\% \\
Innovationsrate & Schnellere Entwicklung neuer Ansätze & Schwer quantifizierbar \\
\bottomrule
\end{tabularx}
\end{table}

\subsubsection{Reduktion von Fehlinvestitionen}

\begin{table}[h]
\centering
\begin{tabular}{ll}
\toprule
\textbf{Ohne BEATRIX} & \textbf{Mit BEATRIX} \\
\midrule
70\% aller Change-Projekte scheitern & Scheiterquote reduziert auf 40--50\% \\
Durchschnittliche Projektkosten: CHF 500k--5M & Frühzeitige Identifikation von Risiken \\
Verschwendung pro Jahr (CH): CHF 1--5 Mrd. & Eingesparte Kosten (CH): CHF 500 Mio.--2 Mrd./Jahr \\
Ursache: Fehlende Verhaltensanalyse & ROI für Volkswirtschaft: Signifikant positiv \\
\bottomrule
\end{tabular}
\end{table}

\subsubsection{Gesellschaftlicher Nutzen}

\begin{table}[h]
\centering
\begin{tabular}{lll}
\toprule
\textbf{Bereich} & \textbf{Aktuell} & \textbf{Mit BEATRIX} \\
\midrule
Policy-Qualität & Intuitionsbasiert & Evidenzbasiert \\
Öffentliche Gesundheit & Geringe Compliance & Verbessertes Verhalten \\
Nachhaltigkeit & Appelbasiert & Verhaltensgesteuert \\
Finanzielle Bildung & Niedrig & Verbesserte Entscheidungen \\
\bottomrule
\end{tabular}
\end{table}

\subsection{Kann die Schweiz in diesem Bereich Leader werden?}

\subsubsection{Die einzigartige Schweizer Position}

\begin{table}[h]
\centering
\begin{tabularx}{\textwidth}{lXX}
\toprule
\textbf{Faktor} & \textbf{Schweiz-Vorteil} & \textbf{Relevanz für BEATRIX} \\
\midrule
Ernst Fehr & Weltführender Verhaltensökonom & Wissenschaftliche Legitimation \\
UZH-Forschung & Top-Rankings global & Methodische Exzellenz \\
FehrAdvice & 15+ Jahre Praxiserfahrung & Umsetzungskompetenz \\
Neutralität & Vertrauenswürdiger Standort & Datenschutz, Reputation \\
Wohlstand & Hohe Beratungsnachfrage & Lokaler Testmarkt \\
Mehrsprachigkeit & DE/FR/IT/EN & Internationale Skalierung \\
\bottomrule
\end{tabularx}
\end{table}

\subsubsection{Wettbewerbsanalyse: Schweiz vs. Alternativen}

\begin{table}[h]
\centering
\begin{tabularx}{\textwidth}{lXXX}
\toprule
\textbf{Land/Region} & \textbf{Stärke} & \textbf{Schwäche} & \textbf{CH-Vorteil} \\
\midrule
USA & Tech-Ökosystem, VC & Datenschutz-Bedenken & \checkmark~Vertrauenswürdigkeit \\
UK (BIT) & Nudge-Pionier & Brexit-Unsicherheit & \checkmark~Stabilität \\
Deutschland & Grosser Markt & Weniger agil & \checkmark~KMU-Geschwindigkeit \\
Singapur & Gov-Tech-Fokus & Klein, kein Fehr & \checkmark~Wissenschaftliche Basis \\
\bottomrule
\end{tabularx}
\end{table}

\subsubsection{Vision: Schweiz als globaler Hub für Behavioral AI}

\begin{table}[h]
\centering
\begin{tabular}{ll}
\toprule
\textbf{Jahr} & \textbf{Meilenstein} \\
\midrule
\textbf{2026} & Schweizer Pionierrolle: Erstes evidenzbasiertes GDSS aus der Schweiz \\
\textbf{2028} & Europäischer Leader: Erste skalierbare Plattform am Markt \\
\textbf{2030} & Globaler Hub: ``Silicon Valley für Behavioral AI'' \\
\bottomrule
\end{tabular}
\end{table}

\begin{summarybox}[Zusammenfassung: Impact for Switzerland]
\textbf{Wettbewerbsfähigkeit:}
\begin{itemize}
    \item Proprietäre Schweizer Plattform
    \item 30--50 direkte Jobs bis 2030
    \item Neue Berufsprofile (Behavioral Engineer, etc.)
    \item CHF 20--50 Mio. Wertschöpfung bis 2030
\end{itemize}

\textbf{KMU-Uni-Zusammenarbeit:}
\begin{itemize}
    \item Musterbeispiel für Innosuisse-Kooperation
    \item Win-Win: Praxis + Wissenschaft
    \item Langfristiges Forschungsprogramm
\end{itemize}

\textbf{Volkswirtschaft \& Gesellschaft:}
\begin{itemize}
    \item Bessere Entscheidungen (+50\%)
    \item Weniger gescheiterte Projekte (--20--30\%)
    \item Demokratisierung von Verhaltensökonomie
\end{itemize}

\textbf{Schweiz als Leader:}
\begin{itemize}
    \item Einzigartige Position (Fehr + UZH + FehrAdvice)
    \item Vision: ``Silicon Valley für Behavioral AI''
\end{itemize}
\end{summarybox}

\newpage

%%%%%%%%%%%%%%%%%%%%%%%%%%%%%%%%%%%%%%%%%%%%%%%%%%%%%%%%%%%%%%%%%%%%%%%%%%%%%%%
\section*{Gesamtzusammenfassung}
\addcontentsline{toc}{section}{Gesamtzusammenfassung}
%%%%%%%%%%%%%%%%%%%%%%%%%%%%%%%%%%%%%%%%%%%%%%%%%%%%%%%%%%%%%%%%%%%%%%%%%%%%%%%

\begin{table}[h]
\centering
\begin{tabularx}{\textwidth}{lX}
\toprule
\textbf{Frage} & \textbf{Kernaussage} \\
\midrule
\textbf{a) Innovation} & BEATRIX transformiert Beratung von Handwerk zu Ingenieurswissenschaft durch axiomatische Verhaltensmodellierung \\
\textbf{b) Markt} & Schliesst Lücke zwischen harten Daten und weichen Faktoren; --85\% Zeit, --60\% Kosten, +45\% Genauigkeit \\
\textbf{c) Machbarkeit} & 20+ Pilotprojekte durchgeführt; PoC validiert Skalierbarkeit (Gall) und Akzeptanz (Luger) \\
\textbf{d) Partnerschaft} & UZH Gall (Technik) + Luger (Management) = ideale Komplementarität; langfristiges Forschungsprogramm \\
\textbf{e) Schweiz-Impact} & Schweiz als Leader in Behavioral AI; neue Jobs; SME-Uni-Zusammenarbeit als Muster \\
\bottomrule
\end{tabularx}
\end{table}

\vspace{1cm}

\begin{center}
\fbox{\parbox{0.9\textwidth}{
\centering
\textbf{Strategischer Kernsatz}\\[0.5em]
\textit{``Dieses Projekt ist keine blosse IT-Entwicklung, sondern liefert die Grundlagen für eine wissenschaftlich fundierte Transformation der Management-Praxis, die nur durch den Einsatz von High-End-Technologie (BEATRIX) möglich wird.''}
}}
\end{center}

\vspace{2cm}

\begin{flushright}
\textit{Dokument erstellt: 18. Januar 2026}\\
\textit{Für: Innosuisse Innovation Cheque}\\
\textit{Fachgebiet: Social Sciences \& Business Management (SSBM)}
\end{flushright}

\end{document}
